\section*{Limitations}
Our study faced several significant limitations that influenced both our approach and the findings. One of the primary constraints was using a 4-bit quantized model due to resource constraints, which restricted our ability to leverage larger parametric models, such as those with 70B or 40B parameters. The token limitation and high subscription charges for paid cloud services further exacerbated this issue, limiting our capacity to perform inference and fine-tuning on more advanced models. This limitation likely restricted the full exploration of these models' capabilities, potentially affecting the depth and quality of the insights and performance metrics we could achieve.

Additionally, the resource-intensive nature of obtaining legal expert annotations presented another challenge. The high costs and significant time required for acquiring these annotations made obtaining expert evaluations for the entire PredEx test dataset impractical. As a result, we use the same 50 random documents as used in \cite{nigam2024legaljudgmentreimaginedpredex} for expert review and Likert score evaluations. While necessary, this approach potentially limits the breadth and depth of our expert evaluation, as it does not encompass the entire dataset.

The applicability of LLMs in the legal domain, particularly for tasks involving legal judgment prediction and explanation, remains uncertain based on our findings. While they show proficiency in conversational contexts, their performance in tasks requiring complex logic or specialized knowledge, such as legal reasoning, is less convincing. Analyzing lengthy legal documents and generating predictions and explanations proved challenging for generative models. This challenge is particularly evident when the models must process and understand intricate legal reasoning and contexts.

Lastly, the dataset used in this study comprised only English-language judgments, excluding other regional languages such as Hindi and Bengali. This limitation underscores the need for more inclusive datasets representing the linguistic diversity in legal documents across different jurisdictions.

These limitations highlight the complexities and challenges of applying LLMs to specialized tasks like legal judgment prediction and explanation. They also underscore the need for ongoing research and development to comprehensively enhance AI models' capabilities in interpreting and understanding legal documents and contexts.
